
\paragraph{\textbf{Remarks 4.6.}}
\SGARefTarget{sga3:xiv:4.6}
Unfortunately, it is no longer true in general that two Borel
subgroups of \(G\) are locally conjugate for the fpqc topology.
This can fail even when \(G\) is affine over \(S\) and \(S\) is
the spectrum of the algebra of dual numbers over an algebraically
closed field \(k\).

As a consequence of this regrettable fact, let us point out that
if \(G\) is a smooth, affine, connected algebraic group over a
nonperfect field \(k\), it is in general impossible to define
naturally a homogeneous space \(D\) under \(G\) that plays the
role of a flag variety, that is, of the variety of Borel subgroups
of \(G\). Over the algebraic closure \(\bar k\) of \(k\), such a
space would have to be isomorphic to \(G_{\bar k}/B\), where \(B\)
is a Borel subgroup of \(G_{\bar k}\).

Indeed, suppose that the quotient of \(G_{\bar k}\) by its radical
\(\bar R\) is an adjoint semisimple group. The kernel of
\[
  G_{\bar k}\longrightarrow
  \operatorname{Aut}_{\bar k}(\bar D)
\]
is then the radical of \(G_{\bar k}\). Consequently, if \(\bar D\)
came from a homogeneous space \(D\) under \(G\), the radical
\(\bar R\) would come from a subgroup \(R\) of \(G\). But examples
are easily constructed for which \(G_{\bar k}/\bar R\) is adjoint
while \(\bar R\) is not defined over \(k\).

Under these conditions, the functor
\[
  \mathfrak B\colon(\operatorname{Sch}/S)^\circ
    \longrightarrow \operatorname{Ens},
  \qquad
  \mathfrak B(S')=
  \{\text{Borel subgroups of }G_{S'}\},
\]
is not representable by a smooth \(S\)-prescheme. From the
infinitesimal point of view \textup{(III, \S3)}, the failure of the
conjugacy theorem is expressed by the fact that, if \(B\) is a
Borel subgroup of the smooth algebraic group \(G\), the cohomology
group
\[
  H^1(B,\mathfrak g/\mathfrak b)
\]
may be nonzero.\footnote{Here \(\mathfrak b\) denotes the Lie
algebra of \(B\).}

We shall see, on the other hand, in a later exposé that when \(G\)
is semisimple, or more generally reductive, such unpleasant
phenomena do not occur. These phenomena, together no doubt with
the absence of good existence theorems, explain why Borel
subgroups play only a relatively unobtrusive role in the
scheme-theoretic study of general group schemes, whereas they
dominate the theory of semisimple group schemes in the later
exposés.

\paragraph{\textbf{Proposition 4.8.}}
\SGARefTarget{sga3:xiv:4.8}
\footnote{There is no item numbered 4.7 in the source.}
{\itshape Let \(G\) be a smooth \(S\)-prescheme in groups of finite
presentation with connected fibers, and let \(B\) be a Borel
subgroup of \(G\). Then \(B\) is its own normalizer and is a closed
subprescheme of \(G\).}

By \SGARefLink{sga3:xii:7.10}{XII~7.10}, it is enough to prove that, over an algebraically
closed field \(k\), every element of \(G(k)\) that normalizes \(B\)
belongs to \(B(k)\). When \(G\) is affine this is a fundamental
result of Chevalley
\textup{\emph{(Bible)}, 9, \SGARefLink{sga3:xiv:7:lemma1}{Theorem~1}}; the general case reduces
immediately to it by the reduction already used in \SGARefLink{sga3:xiv:4.2}{4.2}.

\paragraph{\textbf{Remarks 4.8.1.}}
\SGARefTarget{sga3:xiv:4.8.1}
\SGARefLink{sga3:xiv:4.5}{Definition~4.5} may be generalized by also introducing the notion
of a \emph{parabolic subgroup} of \(G\). This means a subprescheme
in groups \(P\) of \(G\), smooth and of finite presentation over
\(S\), such that for every \(s\in S\) the geometric fiber
\(P_{\bar s}\) is a parabolic subgroup of \(G_{\bar s}\), that is,
contains a Borel subgroup of \(G_{\bar s}\).

\SGARefLink{sga3:xiv:4.8}{Proposition~4.8} extends to a parabolic subgroup \(P\) of \(G\), by
the same reduction to the known set-theoretic assertion. We note
the following consequence \textup{(compare Exposé XVI)}. \SGAInlineRefTarget{sga3:xiv:paragraph:4-8-bis}If \(P\)
is a parabolic subgroup of \(G\), then \(G/P\) is representable by
a quasi-projective prescheme of finite presentation over \(S\);
here \(G\) is assumed to have connected fibers. Moreover, \(G/P\)
is plainly smooth over \(S\), and its geometric fibers are
connected and proper. It follows readily from EGA~III, 5.5.1 that
\[
  D=G/P
\]
is proper, hence projective, over \(S\). If its relative dimension
is \(n\), it is known that the inverse of the invertible sheaf
\(\Omega^n_{D/S}\) induces ample sheaves on the geometric fibers of
\(D/S\); hence, by EGA~III, 4.7.1,
\[
  \bigl(\Omega^n_{D/S}\bigr)^{-1}
\]
is ample on \(D\) relative to \(S\).
